\section{Матроиды. Максимальные независимые подмножества, базисы матроида. Примеры матроидов. Жадный алгоритм.}

\subsection*{Определение матроида}

\paragraph{Определение}
\textbf{Матроидом} называется пара $M = (E, \mathcal{I})$, где $E$ — конечное множество носителя ($|E|=n$), а $\mathcal{I} \subseteq 2^E$ — семейство его подмножеств, называемых \textbf{независимыми}, для которых выполняются три аксиомы:
\begin{itemize}
    \item \textbf{M1:} $\emptyset \in \mathcal{I}$ (пустое множество независимо).
    \item \textbf{M2:} Если $A \in \mathcal{I}$ и $B \subset A$, то $B \in \mathcal{I}$ (подмножество независимого множества независимо).
    \item \textbf{M3 (Аксиома замены):} Если $A, B \in \mathcal{I}$ и $|B| = |A| + 1$, то существует элемент $e \in B \setminus A$, такой что $A \cup \{e\} \in \mathcal{I}$.
\end{itemize}

\paragraph{Пример}
Семейство линейно независимых векторов любого векторного пространства является матроидом. Это исторический прообраз понятия матроида.

\subsection*{Максимальные независимые подмножества и базисы}

\paragraph{Определение}
Пусть $X \subseteq E$. \textbf{Максимальным независимым подмножеством} множества $X$ называется такое $Y \subseteq X$, что $Y \in \mathcal{I}$, и не существует другого $Z \in \mathcal{I}$, такого что $Y \subset Z \subseteq X$. Множество таких подмножеств обозначается $X'$.

\paragraph{Аксиома M4 и равномощность}
Важнейшее свойство матроида: во всяком подмножестве $X$ все максимальные независимые подмножества имеют одинаковую мощность (\textbf{аксиома M4}).
$$\forall X \subseteq E \ (\forall Y, Z \in X' \implies |Y| = |Z|).$$
Новиков доказывает, что при выполнении M1 и M2 аксиомы M3 и M4 эквивалентны.

\paragraph{Базисы}
Максимальные независимые подмножества всего носителя $E$ называются \textbf{базисами} матроида. Согласно свойству M4, все базисы матроида имеют одинаковую мощность (ранг матроида).

\subsection*{Жадный алгоритм}

\paragraph{Постановка задачи}
Задана весовая функция $w: E \to \mathbb{R}^+$. Нужно найти независимое множество $X \in \mathcal{I}$ с максимальным суммарным весом $w(X) = \sum_{e \in X} w(e)$.

\paragraph{Алгоритм}
1. Упорядочить элементы $E$ по убыванию весов: $w(e_1) \ge w(e_2) \ge \dots \ge w(e_n)$.
2. Инициализировать $X = \emptyset$.
3. Для $i$ от $1$ до $n$: если $X \cup \{e_i\} \in \mathcal{I}$, то включить $e_i$ в $X$.

\paragraph{Теорема о жадном алгоритме}
Жадный алгоритм находит оптимальное решение (множество максимального веса) для любой весовой функции $w$ \textbf{тогда и только тогда}, когда структура $(E, \mathcal{I})$ является матроидом.

\subsection*{Примеры и контрпримеры}

\begin{enumerate}
    \item \textbf{Матричный матроид:} $E$ — столбцы матрицы, $\mathcal{I}$ — линейно независимые наборы столбцов. Жадный алгоритм здесь находит базис с максимальным весом.
    \item \textbf{Пример "Один элемент из столбца":} Если задача — выбрать по одному элементу из каждого столбца матрицы с макс. суммой, то это матроид (жадный алгоритм работает).
    \item \textbf{Контрпример "Один из столбца и строки":} Если нужно выбрать элементы так, чтобы из каждой строки и каждого столбца было не более одного элемента, то жадный алгоритм может выдать неоптимальный результат.
\end{enumerate}

\paragraph{Значение}
Матроиды позволяют сразу определить, применим ли эффективный жадный алгоритм ($O(n)$) для решения сложной экстремальной задачи, или же структура задачи требует более тяжелых методов (динамическое программирование, перебор).